\documentclass[12pt]{ctexart}
\usepackage{amsmath, amssymb}
\usepackage{tikz}
\usepackage{geometry}
\usepackage{booktabs}
\usepackage{enumitem}
\geometry{margin=1in}
\usetikzlibrary{arrows.meta}

\newcommand{\overarrow}[1]{\overrightarrow{#1}}
\newcommand{\colvec}[2]{\begin{pmatrix} #1 \\ #2 \end{pmatrix}}

\title{Vectors, Translations, and Meaning\\[2pt] \large Worked Solutions}
\author{Extension Sheet}
\date{}

\begin{document}
\maketitle

% ------------------------------------------------------------------
\section*{1\quad Vector translation review}
% ------------------------------------------------------------------

Translating a point by \(\colvec{a}{b}\) means: add \(a\) to the \(x\)-coordinate and \(b\) to the \(y\)-coordinate.

\begin{enumerate}[label=\textbf{\arabic*.}]
\item For \(A=(2,3)\) translated by \(\mathbf{t}=\colvec{4}{1}\),
\[
A'=(2,3)+\colvec{4}{1}
=(2+4,\;3+1)
=(6,4).
\]
So the destination is \(\boxed{(6,4)}\).

\item For \(B=(-1,5)\) translated by the same vector,
\[
B'=(-1,5)+\colvec{4}{1}
=(-1+4,\;5+1)
=(3,6).
\]
So the destination is \(\boxed{(3,6)}\).
\end{enumerate}

% ------------------------------------------------------------------
\section*{2\quad Word embeddings}
% ------------------------------------------------------------------

Here \(\overarrow{X\ Y}\) means the vector from the first word to the second word, so we compute it as \(Y-X\).

From the grid:
\[
\text{man}=(2,5),\quad
\text{woman}=(3,8),\quad
\text{king}=(7,6),\quad
\text{queen}=(8,9),
\]
\[
\text{prince}=(5,6),\quad
\text{duke}=(4,0),\quad
\text{duchess}=(5,3).
\]

\begin{enumerate}[label=\textbf{\arabic*.}]
\item \[
\overarrow{\text{king}\quad\text{queen}}
=\colvec{8}{9}-\colvec{7}{6}
=\colvec{1}{3}.
\]

\item \[
\overarrow{\text{duke}\quad\text{duchess}}
=\colvec{5}{3}-\colvec{4}{0}
=\colvec{1}{3}.
\]

\item The vector \(\colvec{1}{3}\) is the same in both cases:
\[
\text{king}\to\text{queen}=\colvec{1}{3}
\qquad\text{and}\qquad
\text{duke}\to\text{duchess}=\colvec{1}{3}.
\]
It also maps \text{man} to \text{woman}. Therefore, in this embedding, \(\colvec{1}{3}\) represents the \textbf{male-to-female} direction, i.e.\ a feminine/gender-change vector.

\item Since
\[
\overarrow{\text{man}\quad\text{woman}}
=\colvec{3}{8}-\colvec{2}{5}
=\colvec{1}{3},
\]
yes, it matches the answer to Question 3 exactly.

\item The missing word is at
\[
(5,6)+\colvec{1}{3}=(6,9).
\]
Since \((5,6)\) is \text{prince}, applying the male-to-female vector gives
\[
\text{prince}+\colvec{1}{3}=\text{princess}.
\]
So the missing word is \(\boxed{\text{princess}}\).
\end{enumerate}

% ------------------------------------------------------------------
\section*{3\quad Translating English to Chinese}
% ------------------------------------------------------------------

First read the known coordinates from the grid.

English coordinates:
\[
\text{man}=(9,8),\quad
\text{woman}=(10,11),\quad
\text{boy}=(7,8),\quad
\text{girl}=(8,11),
\]
\[
\text{king}=(8,6),\quad
\text{queen}=(9,9),\quad
\text{prince}=(6,6),\quad
\text{princess}=(7,9).
\]

Chinese coordinates:
\[
\text{男人}=(-6,-2),\quad
\text{女人}=(-5,1),\quad
\text{男孩}=(-8,-2),\quad
\text{女孩}=(-7,1),
\]
\[
\text{国王}=(-7,-4),\quad
\text{王妃}=(-6,-1),\quad
\text{王子}=(-9,-4),\quad
\text{公主}=(-9,0).
\]

\begin{enumerate}[label=\textbf{\arabic*.}]
\item Using the matching pair
\[
\text{man}\ \to\ \text{男人},
\]
the translation vector from English to Chinese is
\[
\mathbf{t}
=\colvec{-6}{-2}-\colvec{9}{8}
=\boxed{\colvec{-15}{-10}}.
\]
Check with another pair:
\[
\text{woman}+\mathbf{t}
=\colvec{10}{11}+\colvec{-15}{-10}
=\colvec{-5}{1},
\]
which is indeed 女人.

\item The completed table is:

\begin{center}
\renewcommand{\arraystretch}{1.3}
\begin{tabular}{@{}llll@{}}
\toprule
English & English coords & Chinese & Chinese coords \\
\midrule
man     & (9,8)    & 男人   & (-6,-2) \\
woman   & (10,11)  & 女人   & (-5,1) \\
boy     & (7,8)    & 男孩   & (-8,-2) \\
girl    & (8,11)   & 女孩   & (-7,1) \\
king    & (8,6)    & 国王   & (-7,-4) \\
queen   & (9,9)    & 王妃   & (-6,-1) \\
prince  & (6,6)    & 王子   & (-9,-4) \\
princess& (7,9)    & 公主   & (-9,0) \\
\bottomrule
\end{tabular}
\end{center}

So the missing labels on the coordinate grid are:
\[
(10,11)=\text{woman},\quad
(9,9)=\text{queen},\quad
(6,6)=\text{prince},\quad
(7,9)=\text{princess},
\]
and on the Chinese side:
\[
(-5,1)=\text{女人},\quad
(-8,-2)=\text{男孩},\quad
(-9,0)=\text{公主}.
\]

\item The vector \(\colvec{-2}{0}\) appears between these English pairs:
\[
\text{man}\to\text{boy},\qquad
\text{woman}\to\text{girl},\qquad
\text{king}\to\text{prince},\qquad
\text{queen}\to\text{princess}.
\]
In each case, the \(x\)-coordinate decreases by \(2\) and the \(y\)-coordinate stays the same. Thus \(\colvec{-2}{0}\) represents the meaning \textbf{“adult/older” to “younger/child”}, or simply the \textbf{age/youth direction} in this embedding.

\item The bottom row does not fit the pattern because English \text{princess} is the feminine form of \text{prince}, but Chinese 公主 is not a regular feminine form built from 王子.

Applying the translation vector to \text{princess} gives
\[
\text{princess}+\mathbf{t}
=\colvec{7}{9}+\colvec{-15}{-10}
=\colvec{-8}{-1},
\]
but the actual position of 公主 is
\[
\colvec{-9}{0}.
\]
The difference is
\[
\colvec{-9}{0}-\colvec{-8}{-1}
=\colvec{-1}{1}.
\]

The Chinese characters explain why: words such as 女人, 女孩, 王妃 all contain the female radical 女, but 公主 does not. Chinese has a separate, lexicalised word for “princess” rather than constructing it as a regular feminine version of 王子. Therefore the regular translation-vector pattern breaks for this row.
\end{enumerate}

\end{document}